新的尝试方向：1.将原有的材料的烧结温度继续推向高温，观察有没有新的有趣的变化；2.尝试将葡萄糖\footnote{https://saylordotorg.github.io/text_general-chemistry-principles-patterns-and-applications-v1.0/s17-04-effects-of-temperature-and-pre.html}引入纳米纤维，通过烧结在纳米纤维网络中引入更多的纳米非晶碳，观察对材料性能的影响。

针对2具体来说，首先应当尝试复刻原论文的实验\cite{du2019a}，分别掺入1mol/L、2mol/L、3mol/L的葡萄糖溶液，观察烧结后材料的性能变化。之后可以尝试更高浓度的葡萄糖溶液，观察是否存在一个最佳浓度范围。
不过需要注意的是，原作者尝试的是高压反应釜路线，我们采用的是常温常压下的静电纺丝路线，因此实验结果可能会有所不同。